% Généré par outils/tables_cli.py — ne pas modifier à la main.
\begin{tabularx}{\linewidth}{@{}>{\ttfamily\small\raggedright}p{0.38\linewidth}X@{}}
\toprule
\multicolumn{2}{@{}p{\linewidth}@{}}{\sffamily\bfseries coupole}\\
\midrule
--version & show the version and exit\\
--lang,\allowbreak{} --langue \{auto,\allowbreak{}fr,\allowbreak{}en\} & message language (auto: system language)\\
-v,\allowbreak{} --verbeux,\allowbreak{} --verbose & detailed log\\
--reinitialiser-interface,\allowbreak{} --reset-interface & erase the saved layout (window, columns, filters, file-dialog folders) and exit; settings do not change\\
\bottomrule
\end{tabularx}
\par\medskip
\begin{tabularx}{\linewidth}{@{}>{\ttfamily\small\raggedright}p{0.38\linewidth}X@{}}
\toprule
\multicolumn{2}{@{}p{\linewidth}@{}}{\sffamily\bfseries coupole astap\newline\normalfont\small Look for ASTAP and its star catalogue, show what to install for this system, or record their location. Coupole works without ASTAP.}\\
\midrule
--guide & show the installation guide for this system\\
--definir,\allowbreak{} --set PATH & path of the ASTAP executable (or its folder)\\
--catalogue,\allowbreak{} --catalog FOLDER & star catalogue folder (D80, D50...)\\
--oublier,\allowbreak{} --forget & forget recorded paths (back to detection)\\
--json & JSON output (for scripts)\\
\bottomrule
\end{tabularx}
\par\medskip
\begin{tabularx}{\linewidth}{@{}>{\ttfamily\small\raggedright}p{0.38\linewidth}X@{}}
\toprule
\multicolumn{2}{@{}p{\linewidth}@{}}{\sffamily\bfseries coupole rapports (reports)\newline\normalfont\small Nothing is sent without consent. Without consent, reports stay in a local folder.}\\
\midrule
--activer,\allowbreak{} --enable & allow sending (anonymous) reports\\
--desactiver,\allowbreak{} --disable & forbid sending reports\\
\bottomrule
\end{tabularx}
\par\medskip
\begin{tabularx}{\linewidth}{@{}>{\ttfamily\small\raggedright}p{0.38\linewidth}X@{}}
\toprule
\multicolumn{2}{@{}p{\linewidth}@{}}{\sffamily\bfseries coupole maj (update)\newline\normalfont\small Query the releases published on GitHub. The standalone package updates itself; a pip/pipx installation shows the command to run.}\\
\midrule
--appliquer,\allowbreak{} --apply & download and install the update (standalone package)\\
\bottomrule
\end{tabularx}
\par\medskip
\begin{tabularx}{\linewidth}{@{}>{\ttfamily\small\raggedright}p{0.38\linewidth}X@{}}
\toprule
\multicolumn{2}{@{}p{\linewidth}@{}}{\sffamily\bfseries coupole sources\newline\normalfont\small Without option: list the addresses and their origin (shipped default, remote file, forced value).}\\
\midrule
--forcer,\allowbreak{} --set KEY VALUE & force a value (never overwritten by an update)\\
--defaut,\allowbreak{} --default KEY & go back to the default value for this key\\
--tout-reinitialiser,\allowbreak{} --reset-all & go back to the default values for every key\\
--tester,\allowbreak{} --test KEY & test the connection to this address\\
--distant,\allowbreak{} --remote & fetch the sources file published in the repository now\\
\bottomrule
\end{tabularx}
\par\medskip
\begin{tabularx}{\linewidth}{@{}>{\ttfamily\small\raggedright}p{0.38\linewidth}X@{}}
\toprule
\multicolumn{2}{@{}p{\linewidth}@{}}{\sffamily\bfseries coupole ohp inventaire (inventory)\newline\normalfont\small Bank inventory from the public TAP service (http://tap-ufe.obspm.fr/tap, table ivoa.obscore), kept in a local cache.}\\
\midrule
--rafraichir,\allowbreak{} --refresh & query the Observatory's TAP service again (one request)\\
--manquantes,\allowbreak{} --missing & list what remains to download compared with the output folder (--dest)\\
--dest FOLDER & destination folder (default: \textasciitilde{}/Coupole/OHP\_DU\_ECU or the one in the settings)\\
--json & JSON output (for scripts)\\
\bottomrule
\end{tabularx}
\par\medskip
\begin{tabularx}{\linewidth}{@{}>{\ttfamily\small\raggedright}p{0.38\linewidth}X@{}}
\toprule
\multicolumn{2}{@{}p{\linewidth}@{}}{\sffamily\bfseries coupole ohp catalogue (catalog)\newline\normalfont\small The 194 database names grouped into objects, sorted by type.}\\
\midrule
--type TYPE & category: ast, neocp, com, pla, tno, pn, neb, amas, gal, eto (repeatable)\\
--telescope \{T120,\allowbreak{}IRIS\} & telescope (repeatable)\\
--chercher,\allowbreak{} --search TEXT & keep objects whose name contains this text\\
--json & JSON output (for scripts)\\
\bottomrule
\end{tabularx}
\par\medskip
\begin{tabularx}{\linewidth}{@{}>{\ttfamily\small\raggedright}p{0.38\linewidth}X@{}}
\toprule
\multicolumn{2}{@{}p{\linewidth}@{}}{\sffamily\bfseries coupole ohp images (list)\newline\normalfont\small One line per image: date, telescope, filter, exposure, flags (duplicate, shared date, daytime).}\\
\midrule
OBJECT & object(s): « NGC 6888 », « Pluto », « 1998 AG6 », or a database name\\
--type TYPE & category: ast, neocp, com, pla, tno, pn, neb, amas, gal, eto (repeatable)\\
--telescope \{T120,\allowbreak{}IRIS\} & telescope (repeatable)\\
--filtre,\allowbreak{} --filter F & filter: B, V, R (Johnson and Cousins), Bc, Vc, Rc, Ha, OIII, g, r, i, z, or the exact name (« Johnson R »); repeatable or comma-separated\\
--nuit,\allowbreak{} --night YYYY-MM-DD & night (evening date); repeatable\\
--sans-dates-douteuses,\allowbreak{} --no-doubtful-dates & leave out exposures with a shared DATE-OBS or dated in daytime\\
--tout,\allowbreak{} --all & the whole bank (≈ 78 GB)\\
--csv FILE & also write the list to this CSV file\\
--dest FOLDER & destination folder (default: \textasciitilde{}/Coupole/OHP\_DU\_ECU or the one in the settings)\\
\bottomrule
\end{tabularx}
\par\medskip
\begin{tabularx}{\linewidth}{@{}>{\ttfamily\small\raggedright}p{0.38\linewidth}X@{}}
\toprule
\multicolumn{2}{@{}p{\linewidth}@{}}{\sffamily\bfseries coupole ohp estimer (estimate)\newline\normalfont\small Estimate before downloading, with the free space at the destination.}\\
\midrule
OBJECT & object(s): « NGC 6888 », « Pluto », « 1998 AG6 », or a database name\\
--type TYPE & category: ast, neocp, com, pla, tno, pn, neb, amas, gal, eto (repeatable)\\
--telescope \{T120,\allowbreak{}IRIS\} & telescope (repeatable)\\
--filtre,\allowbreak{} --filter F & filter: B, V, R (Johnson and Cousins), Bc, Vc, Rc, Ha, OIII, g, r, i, z, or the exact name (« Johnson R »); repeatable or comma-separated\\
--nuit,\allowbreak{} --night YYYY-MM-DD & night (evening date); repeatable\\
--sans-dates-douteuses,\allowbreak{} --no-doubtful-dates & leave out exposures with a shared DATE-OBS or dated in daytime\\
--tout,\allowbreak{} --all & the whole bank (≈ 78 GB)\\
--format \{xisf,\allowbreak{}xisf16,\allowbreak{}fz,\allowbreak{}fits\} & output format: xisf (PixInsight, default), xisf16 (XISF compatible with N.I.N.A. and Siril: 16-bit integers, zlib), fz (lossless compressed FITS), fits (float32)\\
--dest FOLDER & destination folder (default: \textasciitilde{}/Coupole/OHP\_DU\_ECU or the one in the settings)\\
\bottomrule
\end{tabularx}
\par\medskip
\begin{tabularx}{\linewidth}{@{}>{\ttfamily\small\raggedright}p{0.38\linewidth}X@{}}
\toprule
\multicolumn{2}{@{}p{\linewidth}@{}}{\sffamily\bfseries coupole ohp traiter (process)\newline\normalfont\small Full processing of a selection, resumable: running the same command again resumes where it stopped.}\\
\midrule
OBJECT & object(s): « NGC 6888 », « Pluto », « 1998 AG6 », or a database name\\
--type TYPE & category: ast, neocp, com, pla, tno, pn, neb, amas, gal, eto (repeatable)\\
--telescope \{T120,\allowbreak{}IRIS\} & telescope (repeatable)\\
--filtre,\allowbreak{} --filter F & filter: B, V, R (Johnson and Cousins), Bc, Vc, Rc, Ha, OIII, g, r, i, z, or the exact name (« Johnson R »); repeatable or comma-separated\\
--nuit,\allowbreak{} --night YYYY-MM-DD & night (evening date); repeatable\\
--sans-dates-douteuses,\allowbreak{} --no-doubtful-dates & leave out exposures with a shared DATE-OBS or dated in daytime\\
--tout,\allowbreak{} --all & the whole bank (≈ 78 GB)\\
--dest FOLDER & destination folder (default: \textasciitilde{}/Coupole/OHP\_DU\_ECU or the one in the settings)\\
--format \{xisf,\allowbreak{}xisf16,\allowbreak{}fz,\allowbreak{}fits\} & output format: xisf (PixInsight, default), xisf16 (XISF compatible with N.I.N.A. and Siril: 16-bit integers, zlib), fz (lossless compressed FITS), fits (float32)\\
--noms,\allowbreak{} --names \{fr,\allowbreak{}en\} & language of folder names, object names and headers (default: interface language)\\
--astap \{auto,\allowbreak{}tous,\allowbreak{}suspectes,\allowbreak{}jamais,\allowbreak{}all,\allowbreak{}suspicious,\allowbreak{}never\} & ASTAP check: auto/all (every image when ASTAP is present), suspicious, never\\
--telechargements,\allowbreak{} --downloads N & simultaneous downloads (automatic by default; at most 4)\\
--conversions N & simultaneous conversions (default: from cores and memory)\\
--econome,\allowbreak{} --eco & economy mode: 1 conversion, 2 downloads\\
--debit,\allowbreak{} --rate MB/s & maximum download rate (default 8 MB/s)\\
--garder-fits,\allowbreak{} --keep-fits & also keep the original FITS (in \_traitement/fits\_origine)\\
--garder-doublons,\allowbreak{} --keep-duplicates & also process duplicates (inventory and identical pixels) instead of leaving them out\\
--oui,\allowbreak{} --yes & accept a selection over 5 GB without confirmation\\
\bottomrule
\end{tabularx}
\par\medskip
\begin{tabularx}{\linewidth}{@{}>{\ttfamily\small\raggedright}p{0.38\linewidth}X@{}}
\toprule
\multicolumn{2}{@{}p{\linewidth}@{}}{\sffamily\bfseries coupole ohp tout (all)\newline\normalfont\small The whole bank (≈ 8,000 images, ≈ 78 GB of FITS → ≈ 30 GB as XISF), with a volume and time estimate before confirmation; running the same command again resumes where it stopped.}\\
\midrule
--dest FOLDER & destination folder (default: \textasciitilde{}/Coupole/OHP\_DU\_ECU or the one in the settings)\\
--format \{xisf,\allowbreak{}xisf16,\allowbreak{}fz,\allowbreak{}fits\} & output format: xisf (PixInsight, default), xisf16 (XISF compatible with N.I.N.A. and Siril: 16-bit integers, zlib), fz (lossless compressed FITS), fits (float32)\\
--noms,\allowbreak{} --names \{fr,\allowbreak{}en\} & language of folder names, object names and headers (default: interface language)\\
--astap \{auto,\allowbreak{}tous,\allowbreak{}suspectes,\allowbreak{}jamais,\allowbreak{}all,\allowbreak{}suspicious,\allowbreak{}never\} & ASTAP check: auto/all (every image when ASTAP is present), suspicious, never\\
--telechargements,\allowbreak{} --downloads N & simultaneous downloads (automatic by default; at most 4)\\
--conversions N & simultaneous conversions (default: from cores and memory)\\
--econome,\allowbreak{} --eco & economy mode: 1 conversion, 2 downloads\\
--debit,\allowbreak{} --rate MB/s & maximum download rate (default 8 MB/s)\\
--garder-fits,\allowbreak{} --keep-fits & also keep the original FITS (in \_traitement/fits\_origine)\\
--garder-doublons,\allowbreak{} --keep-duplicates & also process duplicates (inventory and identical pixels) instead of leaving them out\\
--oui,\allowbreak{} --yes & accept a selection over 5 GB without confirmation\\
\bottomrule
\end{tabularx}
\par\medskip
\begin{tabularx}{\linewidth}{@{}>{\ttfamily\small\raggedright}p{0.38\linewidth}X@{}}
\toprule
\multicolumn{2}{@{}p{\linewidth}@{}}{\sffamily\bfseries coupole ohp nouveautes (new)\newline\normalfont\small Query the TAP service again, compare with the local copy (output folder) and list what is new; --download processes them into the same tree (never without this option).}\\
\midrule
--telecharger,\allowbreak{} --download & download and convert the listed new images\\
--dest FOLDER & destination folder (default: \textasciitilde{}/Coupole/OHP\_DU\_ECU or the one in the settings)\\
--format \{xisf,\allowbreak{}xisf16,\allowbreak{}fz,\allowbreak{}fits\} & output format: xisf (PixInsight, default), xisf16 (XISF compatible with N.I.N.A. and Siril: 16-bit integers, zlib), fz (lossless compressed FITS), fits (float32)\\
--noms,\allowbreak{} --names \{fr,\allowbreak{}en\} & language of folder names, object names and headers (default: interface language)\\
--astap \{auto,\allowbreak{}tous,\allowbreak{}suspectes,\allowbreak{}jamais,\allowbreak{}all,\allowbreak{}suspicious,\allowbreak{}never\} & ASTAP check: auto/all (every image when ASTAP is present), suspicious, never\\
--telechargements,\allowbreak{} --downloads N & simultaneous downloads (automatic by default; at most 4)\\
--conversions N & simultaneous conversions (default: from cores and memory)\\
--econome,\allowbreak{} --eco & economy mode: 1 conversion, 2 downloads\\
--debit,\allowbreak{} --rate MB/s & maximum download rate (default 8 MB/s)\\
--garder-fits,\allowbreak{} --keep-fits & also keep the original FITS (in \_traitement/fits\_origine)\\
--garder-doublons,\allowbreak{} --keep-duplicates & also process duplicates (inventory and identical pixels) instead of leaving them out\\
--oui,\allowbreak{} --yes & accept a selection over 5 GB without confirmation\\
\bottomrule
\end{tabularx}
\par\medskip
\begin{tabularx}{\linewidth}{@{}>{\ttfamily\small\raggedright}p{0.38\linewidth}X@{}}
\toprule
\multicolumn{2}{@{}p{\linewidth}@{}}{\sffamily\bfseries coupole ohp reorganiser (reorganise, reorganize)\newline\normalfont\small Walk a folder, recognise files produced by Coupole (header) and move them into the stack tree of --dest; name conflict → suffix and a log line; nothing is overwritten.}\\
\midrule
FOLDER & folder holding the files to sort\\
--dest FOLDER & destination folder (default: \textasciitilde{}/Coupole/OHP\_DU\_ECU or the one in the settings)\\
\bottomrule
\end{tabularx}
\par\medskip
\begin{tabularx}{\linewidth}{@{}>{\ttfamily\small\raggedright}p{0.38\linewidth}X@{}}
\toprule
\multicolumn{2}{@{}p{\linewidth}@{}}{\sffamily\bfseries coupole ohp telecharger (download)\newline\normalfont\small Layout: <folder>/<object>/<night>\_<telescope>/<original file>.}\\
\midrule
OBJECT & object(s): « NGC 6888 », « Pluto », « 1998 AG6 », or a database name\\
--type TYPE & category: ast, neocp, com, pla, tno, pn, neb, amas, gal, eto (repeatable)\\
--telescope \{T120,\allowbreak{}IRIS\} & telescope (repeatable)\\
--filtre,\allowbreak{} --filter F & filter: B, V, R (Johnson and Cousins), Bc, Vc, Rc, Ha, OIII, g, r, i, z, or the exact name (« Johnson R »); repeatable or comma-separated\\
--nuit,\allowbreak{} --night YYYY-MM-DD & night (evening date); repeatable\\
--sans-dates-douteuses,\allowbreak{} --no-doubtful-dates & leave out exposures with a shared DATE-OBS or dated in daytime\\
--tout,\allowbreak{} --all & the whole bank (≈ 78 GB)\\
--dest FOLDER & destination folder (default: \textasciitilde{}/Coupole/OHP\_DU\_ECU or the one in the settings)\\
--telechargements,\allowbreak{} --downloads N & simultaneous downloads (automatic by default; at most 4)\\
--debit,\allowbreak{} --rate MB/s & maximum download rate (default 8 MB/s)\\
--oui,\allowbreak{} --yes & accept a selection over 5 GB without confirmation\\
\bottomrule
\end{tabularx}
\par\medskip
\begin{tabularx}{\linewidth}{@{}>{\ttfamily\small\raggedright}p{0.38\linewidth}X@{}}
\toprule
\multicolumn{2}{@{}p{\linewidth}@{}}{\sffamily\bfseries coupole ohp anomalies (anomalies-report)\newline\normalfont\small Duplicates (same file, daytime copy, identical pixels), shared or daytime dates, inconsistent field, classification to check, unknown instrument. Nothing is deleted.}\\
\midrule
--dest FOLDER & add the pixel duplicates found in this processing folder\\
--csv FILE & also write the list to this CSV file\\
--nouveaux,\allowbreak{} --new & only images new since the last refresh\\
\bottomrule
\end{tabularx}
\par\medskip
\begin{tabularx}{\linewidth}{@{}>{\ttfamily\small\raggedright}p{0.38\linewidth}X@{}}
\toprule
\multicolumn{2}{@{}p{\linewidth}@{}}{\sffamily\bfseries coupole ohp classer (classify)\newline\normalfont\small Without argument: list the names to check. With a database name and --object: record the correction (merge into an existing object if --type is omitted). --online: query JPL SBDB and CDS Sesame.}\\
\midrule
DB\_NAME & target name as it appears in the database\\
--objet,\allowbreak{} --object OBJECT & canonical object to associate (existing one: merge)\\
--type \{ast,\allowbreak{}neocp,\allowbreak{}com,\allowbreak{}pla,\allowbreak{}tno,\allowbreak{}pn,\allowbreak{}neb,\allowbreak{}amas,\allowbreak{}gal,\allowbreak{}eto,\allowbreak{}autre\} & category: ast, neocp, com, pla, tno, pn, neb, amas, gal, eto (repeatable)\\
--oublier,\allowbreak{} --forget & forget the correction of this name\\
--en-ligne,\allowbreak{} --online & resolve unknown names online (JPL SBDB, CDS Sesame), with cache\\
\bottomrule
\end{tabularx}
\par\medskip
\begin{tabularx}{\linewidth}{@{}>{\ttfamily\small\raggedright}p{0.38\linewidth}X@{}}
\toprule
\multicolumn{2}{@{}p{\linewidth}@{}}{\sffamily\bfseries coupole ohp ranger (sort)\newline\normalfont\small Useful after an interruption, or to merge several runs in the same folder.}\\
\midrule
--dest FOLDER & destination folder (default: \textasciitilde{}/Coupole/OHP\_DU\_ECU or the one in the settings)\\
\bottomrule
\end{tabularx}
\par\medskip
\begin{tabularx}{\linewidth}{@{}>{\ttfamily\small\raggedright}p{0.38\linewidth}X@{}}
\toprule
\multicolumn{2}{@{}p{\linewidth}@{}}{\sffamily\bfseries coupole ohp bilan (status)\newline\normalfont\small progress of a processing folder}\\
\midrule
--dest FOLDER & destination folder (default: \textasciitilde{}/Coupole/OHP\_DU\_ECU or the one in the settings)\\
--json & JSON output (for scripts)\\
\bottomrule
\end{tabularx}
\par\medskip
\begin{tabularx}{\linewidth}{@{}>{\ttfamily\small\raggedright}p{0.38\linewidth}X@{}}
\toprule
\multicolumn{2}{@{}p{\linewidth}@{}}{\sffamily\bfseries coupole ohp fusionner (merge)\newline\normalfont\small merge the state database on the share with the local working database (after a divergence)}\\
\midrule
--dest FOLDER & destination folder (default: \textasciitilde{}/Coupole/OHP\_DU\_ECU or the one in the settings)\\
\bottomrule
\end{tabularx}
\par\medskip
\begin{tabularx}{\linewidth}{@{}>{\ttfamily\small\raggedright}p{0.38\linewidth}X@{}}
\toprule
\multicolumn{2}{@{}p{\linewidth}@{}}{\sffamily\bfseries coupole ohp metadonnees (metadata)\newline\normalfont\small For each XISF/FITS of the folder: FOCALLEN matched to the measured scale (old value in HISTORY) and the XISF properties Instrument:Sensor:XPixelSize, Instrument:Camera:XBinning, Instrument:Telescope:FocalLength… Pixels are not touched (read back and compared), atomic writing, log in \_traitement/metadonnees.csv. Without --rewrite: dry run.}\\
\midrule
FOLDER & output folder (or one stack)\\
--reecrire,\allowbreak{} --rewrite & really write (otherwise: only say what would change)\\
--noms,\allowbreak{} --names \{fr,\allowbreak{}en\} & language of folder names, object names and headers (default: interface language)\\
\bottomrule
\end{tabularx}
\par\medskip
\begin{tabularx}{\linewidth}{@{}>{\ttfamily\small\raggedright}p{0.38\linewidth}X@{}}
\toprule
\multicolumn{2}{@{}p{\linewidth}@{}}{\sffamily\bfseries coupole qualite\newline\normalfont\small Optional: background, noise, FWHM and ellipticity (3×3 map), background gradient, saturation, trails, sampling}\\
\midrule
FOLDER & folder of stacks (or one image file)\\
--ecrire,\allowbreak{} --write & write QUALITE.csv and QUALITE.txt in each stack\\
--echantillon,\allowbreak{} --sample N & measure N images per stack (spread in time) instead of all\\
--tout,\allowbreak{} --all & measure everything even above 200 images (otherwise: sample of 5 per stack, announced)\\
--processus,\allowbreak{} --processes N & number of measuring processes (default: machine plan)\\
\bottomrule
\end{tabularx}
\par\medskip
\begin{tabularx}{\linewidth}{@{}>{\ttfamily\small\raggedright}p{0.38\linewidth}X@{}}
\toprule
\multicolumn{2}{@{}p{\linewidth}@{}}{\sffamily\bfseries coupole donnees lire (read)\newline\normalfont\small describe the contents of a file (FITS, CSV)}\\
\midrule
FILE & file to read\\
--json & JSON output (for scripts)\\
\bottomrule
\end{tabularx}
\par\medskip
\begin{tabularx}{\linewidth}{@{}>{\ttfamily\small\raggedright}p{0.38\linewidth}X@{}}
\toprule
\multicolumn{2}{@{}p{\linewidth}@{}}{\sffamily\bfseries coupole donnees exporter (export)\newline\normalfont\small export a spectrum or a series as CSV}\\
\midrule
FILE & file to read\\
--csv FILE & export a spectrum or a series as CSV\\
--hdu HDU & index of the dataset in the file (default: the first 1D one)\\
--vitesse,\allowbreak{} --velocity & add the radio velocity (frequency spectrum)\\
--f0 MHz & rest frequency in MHz (default: RESTFRQ, else H I)\\
\bottomrule
\end{tabularx}
\par\medskip
\begin{tabularx}{\linewidth}{@{}>{\ttfamily\small\raggedright}p{0.38\linewidth}X@{}}
\toprule
\multicolumn{2}{@{}p{\linewidth}@{}}{\sffamily\bfseries coupole cosmo\newline\normalfont\small From redshift to distances: comoving, luminosity, angular, lookback time, age, volume, distance modulus, scale (Planck 2018 and other parameter sets)}\\
\midrule
Z & redshift(s) to compute\\
--objet,\allowbreak{} --object NAME & object name: redshift requested from SIMBAD\\
--modele,\allowbreak{} --model \{planck18,\allowbreak{}planck15,\allowbreak{}wmap9,\allowbreak{}simple,\allowbreak{}perso\} & parameter set (default: planck18)\\
--h0 H0 & H₀ in km s⁻¹ Mpc⁻¹ (custom model)\\
--om OM & total Ωm (custom model)\\
--ok OK & Ωk, curvature (planck18 or perso)\\
--shoes & add the SH0ES comparison (planck18)\\
--sans-incertitudes,\allowbreak{} --no-uncertainties & without uncertainties (faster)\\
--csv FILE & write the results to this CSV file\\
--courbes,\allowbreak{} --curves ZMIN ZMAX N & write the curves (z from ZMIN to ZMAX, N points) to the CSV file given by --csv\\
--json & JSON output (for scripts)\\
\bottomrule
\end{tabularx}
\par\medskip
\begin{tabularx}{\linewidth}{@{}>{\ttfamily\small\raggedright}p{0.38\linewidth}X@{}}
\toprule
\multicolumn{2}{@{}p{\linewidth}@{}}{\sffamily\bfseries coupole sites\newline\normalfont\small Observing sites (OpenStreetMap map), time zones, UTC and local time}\\
\midrule
--ajouter,\allowbreak{} --add X & add a site: ID NAME LAT LON ALT TIMEZONE [MPC]\\
--supprimer,\allowbreak{} --delete ID & delete an added site\\
--heure,\allowbreak{} --time ISO\_DATE & show a UTC date (ISO) in site time and evening date\\
--site ID & site identifier (default: ohp)\\
\bottomrule
\end{tabularx}
\par\medskip
\begin{tabularx}{\linewidth}{@{}>{\ttfamily\small\raggedright}p{0.38\linewidth}X@{}}
\toprule
\multicolumn{2}{@{}p{\linewidth}@{}}{\sffamily\bfseries coupole machine\newline\normalfont\small Hardware diagnosis, chosen parallelism, graphics card, ASTAP}\\
\midrule
--json & JSON output (for scripts)\\
\bottomrule
\end{tabularx}
\par\medskip
